% Work/MAT | 1. rocnik | 21. hodina
% Novy navrh: 20 min prubezny test z mnozin, 25 min ciste motivacni diskuse.
% Pokracovani: 22 definice; 23 presne hodnoty; 24 aplikace.
% Ucebnice: ucebnice/04_Goniometrie_a_trigonometrie/Goniometricke funkce.tex.
% Metodika: Realisticky.cz 010119, 040201; zadani a scenar jsou vlastni.
\documentclass[11pt,a4paper]{article}
% Spolecny styl pro ucitelske pripravy MAT (Work/MAT)
% Verze 1.1, 2026-10-23
\NeedsTeXFormat{LaTeX2e}
\ProvidesPackage{priprava}[2026/10/23 v1.1 Sablona ucitelskych priprav MAT]

% Zakladni typografie a cestina
\RequirePackage[T1]{fontenc}
\RequirePackage[utf8]{inputenc}
\RequirePackage[czech]{babel}
\RequirePackage{lmodern}
\RequirePackage{microtype}
\RequirePackage{amsmath,amssymb}
\RequirePackage{xcolor}
\RequirePackage{enumitem}
\RequirePackage[most]{tcolorbox}
\RequirePackage{tabularx,booktabs,array}
\RequirePackage{fancyhdr}
\RequirePackage{needspace}
\RequirePackage{geometry}
\RequirePackage[hidelinks]{hyperref}

\geometry{a4paper,top=17mm,bottom=18mm,left=23mm,right=23mm}
\definecolor{notegray}{gray}{0.39}
\setlength{\parindent}{0pt}
\setlength{\parskip}{2.5pt plus 1pt minus 0.5pt}
\setlist[itemize]{topsep=2pt,itemsep=1pt,parsep=0pt,leftmargin=1.45em}
\setlist[enumerate]{topsep=2pt,itemsep=2pt,parsep=0pt,leftmargin=1.65em}
\setcounter{secnumdepth}{0}
\renewcommand{\arraystretch}{1.12}

% Zahlavi: vlevo rocnik | globalni tema, vpravo cislo: nazev hodiny
% Zapati zustava zcela prazdne.
\pagestyle{fancy}
\fancyhf{}
\renewcommand{\headrulewidth}{0.35pt}
\renewcommand{\footrulewidth}{0pt}
\setlength{\headheight}{14pt}
\newcommand{\nastavPripravu}[4]{%
  \fancyhead[L]{\small #1\enspace\textbar\enspace #2}%
  \fancyhead[R]{\small #3: #4}%
}

% Minimalisticke cernobile boxy. Kazda definice zvlast.
\tcbset{commonlesson/.style={enhanced,breakable=false,colback=white,colframe=black,
  colbacktitle=white,coltitle=black,boxrule=0.45pt,
  arc=0.8mm,outer arc=0.8mm,left=3mm,right=3mm,top=1.6mm,bottom=1.6mm,
  toptitle=1mm,bottomtitle=0.7mm,fonttitle=\bfseries,
  attach boxed title to top left={xshift=2.5mm,yshift=-1.15mm},
  boxed title style={colback=white,colframe=white,boxrule=0pt,left=1mm,right=1mm}}}
\newtcolorbox{zadaniBox}[1][]{commonlesson,title={Zadání},#1}
\newtcolorbox{prikladBox}[1][]{commonlesson,title={Příklad},#1}
\newtcolorbox{definiceBox}[1][]{commonlesson,title={Definice},#1}

% Spolecny box pro zadani a jeho reseni, bez mezery mezi castmi.
% Pouziti:
%   \begin{prikladReseniBox}
%     Zadani ulohy.
%     \tcblower
%     \begin{align*} ... \end{align*}
%   \end{prikladReseniBox}
% Spodni cast ma jemny sedy podklad a delici linku; lze v ni sazet
% vice radku matematiky. Obe casti maji jen jeden spolecny ramecek.
\tcbset{solutionlesson/.style={
  commonlesson,
  enhanced jigsaw,
  breakable,
  colbacklower=black!3,
  segmentation style={solid,draw=black!25,line width=0.35pt},
  fontlower=\small\color{notegray},
  before lower={\noindent\textbf{Řešení:}\enspace},
  leftlower=3mm,rightlower=3mm
}}
\newtcolorbox{prikladReseniBox}[1][]{solutionlesson,title={Příklad},#1}
\newtcolorbox{zadaniReseniBox}[1][]{solutionlesson,title={Zadání},#1}

% Pomucky pro autora pripravy: sedive poznamky; \reseni pro starsi pripravy
\newcommand{\poznamka}[1]{\par\begingroup\small\color{notegray}\itshape
  \textnormal{\textbf{Poznámka pro učitele.}}\enspace #1\par\endgroup}
\newcommand{\reseni}[1]{\par\begingroup\small\color{notegray}
  \textbf{Očekávané řešení:}\enspace #1\par\endgroup}
\newcommand{\blok}[1]{\needspace{5\baselineskip}\subsection*{#1}}
\newcommand{\vyrok}[1]{\uv{#1}}

\nastavPripravu{1. ročník}{Množiny / goniometrie}{21}{Test a úvodní diskuse}
\begin{document}
\blok{0.--20. minuta: průběžný test z množin}
\begin{zadaniBox}[title={Organizace pro učitele}]
\textbf{20 minut čistého času} na samostatný průběžný test z množin,
Vennových diagramů, intervalů a jejich operací (učivo hodin 12--20).
Použij připravené školní zadání; \textbf{tato příprava test nevytváří}.
Po skončení test vyber. Goniometrie není součástí jeho hodnocení.
\end{zadaniBox}
\poznamka{Před vyučováním nachystej testové varianty a pomůcky. Během testu
vysvětluj pouze organizační pokyny. Rozbor písemky nech na dobu,
až budeš znát skutečné chyby; v druhé části dnes otevřeme nové téma.}

\blok{20.--27. minuta: co už žáci vědí?}
Začni otevřenými otázkami, nikoli definicemi. Odpovědi pouze zapisuj na tabuli;
zatím nerozhoduj, kdo zná „správný vzorec“.
\begin{zadaniBox}[title={Otázky pro společnou diskusi}]
\begin{enumerate}[label=\arabic*)]
\item Kdo se už setkal s pojmy \emph{sinus, kosinus, tangens, kotangens}?
  V matematice, fyzice, dílně, při hrách nebo při programování?
\item Co se podle vás dá zjistit, známe-li nějaký úhel?
\item Jaké úhly v trojúhelníku už bezpečně poznáme? Co říká údaj $90^\circ$?
\item Jak spolu může souviset \emph{sklon} a \emph{délka}?
\end{enumerate}
\end{zadaniBox}
\poznamka{Do diskuse mohou žáci přinést i nepřesné představy.
Neopravuj vše vzorcem. Cílem je odhadnout předchozí zkušenosti,
ne klasifikovat znalost goniometrických funkcí.}

\blok{27.--38. minuta: dvě situace bez výpočtu}
\begin{zadaniBox}[title={Situace A: výška, kterou nelze přímo změřit}]
Stojíme na rovném chodníku před školou. Známe svou vzdálenost od budovy
a dovedeme změřit úhel, pod kterým vidíme její střechu.
\emph{Jak bychom mohli určit výšku budovy bez výstupu na střechu?}
Co bychom museli zakreslit? Které údaje ještě potřebujeme znát?
\end{zadaniBox}
\begin{zadaniBox}[title={Situace B: dvě stejně strmé rampy}]
Krátká rampa překoná výšku $1\,\mathrm m$ na vodorovné délce $4\,\mathrm m$.
Druhá překoná $2\,\mathrm m$ na $8\,\mathrm m$.
Jsou stejně strmé, přestože nejsou stejně velké? Jak to poznáme bez úhloměru?
\end{zadaniBox}
\poznamka{Očekávané náměty: náčrt pravoúhlého trojúhelníku,
vodorovná a svislá délka, výška očí, poměr změn $1:4=2:8$.
Záměrně zatím nepoužívej definice $\sin$, $\cos$, $\tan$, $\cot$
ani tabulku hodnot. Nápad s poměrem předznamená 22. hodinu.}

\blok{38.--45. minuta: otázky na příště}
\begin{zadaniBox}[title={Výstupní lístek, bez známky}]
Každý napiš dvě věty: (a) Kde by goniometrie mohla být užitečná?
(b) Co by ses chtěl o měření úhlů a stran dozvědět?
\end{zadaniBox}
Závěrem ukaž jen \textbf{směr dalších tří hodin}: zítra zavedeme pojmy
a čtyři poměry v pravoúhlém trojúhelníku; poté odvodíme hodnoty
pro známé úhly a naučíme se kalkulačku; nakonec je použijeme v úlohách.
\poznamka{Žádný nový formální matematický pojem dnes nezapisujeme do sešitu
jako definici. Časový plán: test 20, úvodní diskuse 7, situace 11,
výstup a uzavření 7 minut. Domácí výpočty nezadávej.}
\end{document}
